\section{Overall Algorithm}
\label{sec:overall}

The receiver combines the estimators and message updates in
Section~\ref{sec:algorithm} with the trained networks in
Section~\ref{sec:learning}. The phase
increments and noise variance are estimated once per access block,
followed by MP iterations between linear estimation and channel and
activity denoising.

\subsection{Message Scheduling}
\label{sec:schedule}

From the original received pilots, estimate $\widehat\sigma^2$ by
\eqref{eq:noise-mdl}, fit the inter-frame phase increments using
\eqref{eq:phase-variable-projection}, and form
$\widehat{\mathbf A}$. Normalize the observation and estimated noise variance
by $s$ and $s^2$, respectively. Initialize
$\mathbf X^{{\rm pri},(0)}=\mathbf0$ and
$\gamma^{(0)}=\rho\widehat P_H$, where $\widehat P_H$ is the training
average of $\norm{H}_F^2/D$ in normalized units.
The initial sites are $E_k^{(0)}=0$ and $c_k^{(0)}=1$ in these units.
An activity history $w_k^{(0)}=\rho$ is maintained for the stopping
check in Section~\ref{sec:damping}.

At each iteration, the linear estimator produces a Gaussian extrinsic message
using \eqref{eq:linear-mean}--\eqref{eq:linear-fraction}.
For $v\in[v_{\min},v_{\max}]$, each trained network is evaluated once
per terminal. The channel error estimates are averaged according to
\eqref{eq:pooled-variance}, and the sites are replaced using
\eqref{eq:site-update} if $\tau<v$; otherwise, the previous sites
are retained. The classifier uses the current $(R_k,v,F_k)$ in
either case.

For $v$ outside the training range, the sites are retained and the
activity probability is obtained from their Gaussian densities:
\begin{equation}
 \begin{aligned}
 \ell_k^{\rm G}&=
 \frac{\norm{R_k}_F^2}{v}
 -\frac{\norm{R_k-E_k}_F^2}{c_k+v}
 -D\log\left(1+\frac{c_k}{v}\right),\\
 p_k&=\sigmoid(\logit\rho+\ell_k^{\rm G}).
 \end{aligned}
 \label{eq:gaussian-evidence}
\end{equation}
The active density is $\CN(\vect(R_k);\vect(E_k),(c_k+v)\I_D)$,
obtained by integrating the Gaussian site against the Gaussian likelihood.
Its ratio to $q_{0,v}$ gives the approximate activity evidence in
\eqref{eq:gaussian-evidence}.

The joint estimate and feedback message are then calculated from the
conditional moments and activity probabilities. The complete procedure
is summarized in Algorithm~1.

\begin{center}
\begin{minipage}{0.96\linewidth}
\hrule\vspace{5pt}
\textbf{Algorithm 1: Joint activity detection and channel estimation}\par
\textbf{Input:} received pilots $\mathbf Y$, original pilot tensor $P$,
and activity marginal $\rho$;
trained $S_\theta,T_\psi,z_\phi$ and fixed receiver settings.\par
\textbf{Output:} channel estimates, activity probabilities and decisions.\par
\begin{enumerate}[leftmargin=1.8em,itemsep=2pt,topsep=4pt]
 \item Estimate $\widehat\sigma^2$ using \eqref{eq:noise-mdl}.
 Fit $\widehat{\bm\beta}$ by the coordinate sweeps and joint gain
 refits in Section~\ref{sec:frontend}; construct $\widehat{\mathbf A}$.
 \item Normalize $\mathbf Y$ and $\widehat\sigma^2$ by $s$ and $s^2$.
 Initialize the prior, sites, activity history, and stability counters.
 Compute the thin SVD of $\widehat{\mathbf A}$ and hold both estimated
 observation parameters fixed.
 \item For $i=1,\ldots,I_{\max}$, compute the linear extrinsic message using
 \eqref{eq:linear-mean}--\eqref{eq:linear-fraction} and damp it by
 \eqref{eq:input-damping}.
 \item If $v$ is within the training range, evaluate $F_k$, $\widehat\tau_k$, and
 $z_\phi(R_k,v,F_k)$; average the conditional mean-squared error
 estimates over terminals and apply the lower bound in
 \eqref{eq:pooled-variance}. Accept the site update
 \eqref{eq:site-update} when $\tau<v$ and otherwise retain the sites.
 Compute $p_k$ by \eqref{eq:activity-probability}.
 \item If $v$ is outside the training range, retain the sites and compute $p_k$
 by \eqref{eq:gaussian-evidence}.
 \item Compute $\mu_k,t_k$ by \eqref{eq:conditional-moments} and
 $\widehat X_k,\mathbf C_{X,k}$ by
 \eqref{eq:mixture}--\eqref{eq:joint-spatial-covariance}.
 \item Form the outgoing message using
 \eqref{eq:joint-average-variance}--\eqref{eq:nonlinear-extrinsic}
 and apply \eqref{eq:output-damping}.
 \item Update the stability counters and activity history.
 Stop when the criterion in Section~\ref{sec:damping} is satisfied.
 \item[]\textbf{end for}
 \item Return $s\widehat X_k$, $s\mu_k$, $p_k$,
 $\widehat\alpha_k=\mathbf1\{p_k\geq1/2\}$, and the stopping status.
 Rescale covariance outputs by $s^2$.
\end{enumerate}
\vspace{3pt}\hrule
\end{minipage}
\end{center}

\subsection{Damping and Stopping Criteria}
\label{sec:damping}

We apply damping to the means and variances of both the denoiser input
and the feedback message.
With new-message weight $0<\eta\leq1$, the input update for $i\geq2$ is
\begin{equation}
 \begin{aligned}
 R_k^{(i)}&=(1-\eta)R_k^{(i-1)}+\eta R_{{\rm L},k}^{(i)},\\
 v^{(i)}&=(1-\eta)v^{(i-1)}+\eta v_{\rm L}^{(i)}.
 \end{aligned}
 \label{eq:input-damping}
\end{equation}
Similarly, the feedback parameters are updated as
\begin{equation}
 \begin{aligned}
 X_k^{{\rm pri},(i)}
 &=(1-\eta)X_k^{{\rm pri},(i-1)}+\eta X_{k,\rm out}^{(i)},\\
 \gamma^{(i)}
 &=(1-\eta)\gamma^{(i-1)}+\eta\gamma_{\rm out}^{(i)}.
 \end{aligned}
 \label{eq:output-damping}
\end{equation}
The first updates are used without damping.

The stopping rule tracks the estimate, activity probabilities, and
input variance. For $i\geq2$, define
\begin{equation}
 \begin{aligned}
 d_X^{(i)}
 &=\frac{\norm{\widehat X^{(i)}-\widehat X^{(i-1)}}_F}
 {\max\{\norm{\widehat X^{(i)}}_F,
         \norm{\widehat X^{(i-1)}}_F,\epsilon_{\rm num}\}},\\
 d_p^{(i)}&=\max_k|p_k^{(i)}-p_k^{(i-1)}|,\\
 d_v^{(i)}&=\frac{|v^{(i)}-v^{(i-1)}|}
 {\max\{v^{(i)},v^{(i-1)}\}},
 \end{aligned}
 \label{eq:iterate-stopping}
\end{equation}
where $\epsilon_{\rm num}>0$ prevents division by zero.
The estimate-stability counter increases when
$d_X^{(i)}<\epsilon_{\rm est}$,
$d_v^{(i)}<\epsilon_{\rm est}$, and
$d_p^{(i)}<\epsilon_p$, and resets otherwise.

We also check the accepted, undamped site parameters.
Let $\lambda_k=c_k^{-1}$ and $J_k=E_k/c_k$, and define
\begin{equation}
 \begin{aligned}
 d_{\lambda,k}^{(i)}
 &=\frac{|\lambda_k^{(i)}-\lambda_k^{(i-1)}|}
 {\max\{\lambda_k^{(i)},\lambda_k^{(i-1)}\}},\\
 d_{J,k}^{(i)}
 &=\frac{\norm{J_k^{(i)}-J_k^{(i-1)}}_F}
 {\max\{1,\norm{J_k^{(i)}}_F,\norm{J_k^{(i-1)}}_F\}},\\
 d_{\rm site}^{(i)}
 &=\max_k\{d_{\lambda,k}^{(i)},d_{J,k}^{(i)}\},\qquad
 d_w^{(i)}=\max_k|p_k^{(i)}-w_k^{(i-1)}|.
 \end{aligned}
 \label{eq:site-stopping}
\end{equation}
The site-stability counter increases only if the current site update is
accepted, $d_{\rm site}^{(i)}<\epsilon_{\rm site}$, and
$d_w^{(i)}<\epsilon_w$; it resets otherwise.
After this check, update
\begin{equation}
 w_k^{(i)}=(1-\zeta)w_k^{(i-1)}+\zeta p_k^{(i)},
 \qquad 0<\zeta\leq1.
 \label{eq:activity-history}
\end{equation}
This history is used only for stopping; the receiver reports $p_k^{(i)}$
and retains the fixed prior $\rho$ in \eqref{eq:activity-probability}.

Iteration stops after at least $I_{\min}$ iterations when both stability
counters reach $L_s$ and the hard activity vector has remained unchanged
for $L_d$ consecutive iterates. The two stability counters reset
independently. If the iteration limit $I_{\max}$ is reached first, the
final iterate is returned with that status. The damping weights,
tolerances, and window lengths are given in the simulation setup.

\subsection{Computational Complexity Analysis}

We analyze the costs of phase estimation, noise estimation, and the
subsequent MP iterations. Let $G_{\rm c}$ be the number of path matrices on the coarse
delay--Doppler grid, $B\geq Q$ the number of phase bins, and $S$ the
number of coordinate sweeps. With dense path-matrix evaluation,
precomputing these path matrices and their pilot responses costs
$O(G_{\rm c}N^3+KG_{\rm c}QN^2)$. Each user search forms correlations
at cost $O(G_{\rm c}QNM)$ and evaluates phase bins by fast Fourier
transforms (FFTs) at cost $O(G_{\rm c}MB\log B)$. An initial search
orders the users, followed by $S$ sweeps.

Let $E_\phi$ count all path-matrix evaluations for alias ranking and
local refinement across the block, including derivative evaluations
up to a constant factor. Evaluating a path matrix, computing its pilot response,
and fitting the gain cost $O(N^3+QN^2+QNM)$. For $K\leq QN$, each joint
gain refit costs $O(QNK^2+QNKM)$. The resulting complexity of phase estimation is
\begin{equation}
 \begin{aligned}
 C_{\rm phase}=O\big(&G_{\rm c}N^3+KG_{\rm c}QN^2\\
 &+(S+1)KG_{\rm c}M(QN+B\log B)\\
 &+E_\phi(N^3+QN^2+QNM)+S(QNK^2+QNKM)\big).
 \end{aligned}
 \label{eq:phase-complexity}
\end{equation}
The local iteration budget bounds optimizer iterations; the actual
number of objective evaluations also depends on line searches.
Noise estimation costs $O(QNM^2+M^3)$, and the thin SVD of
$\widehat{\mathbf A}$ costs $O(KNQ^2)$ for $Q\leq KN$.

In each MP iteration, the dense linear transforms on $MN$
measurement vectors cost $O(QKMN^2)$. The factors in
\eqref{eq:linear-fraction} reuse the initial SVD.

Let $C_\theta$, $C_\psi$, and $C_\phi$ denote the costs of one
terminal-wise forward pass through the mean, error, and activity
networks. Within the training variance range, each iteration evaluates the
networks once per terminal and reuses $F_\theta$ in the classifier.
Taking into account the spatial covariance products and small-matrix
operations, the computational complexity per iteration is given by
\begin{equation}
 C_{\rm iter}=O\!\left(QKMN^2+KM^2N^2+M^3+
 K(C_\theta+C_\psi+C_\phi)\right).
 \label{eq:iteration-complexity}
\end{equation}
The networks are not evaluated outside the training variance range.
For $I$ MP iterations, the complete online cost is
\begin{equation}
 O\!\left(C_{\rm phase}+QNM^2+M^3+KNQ^2+I C_{\rm iter}\right).
 \label{eq:total-complexity}
\end{equation}
Excluding trained parameters and network activations, MP
requires storage of $O(QKN+KMN^2+KM^2)$ entries. The phase estimator additionally
caches $O(G_{\rm c}N^2+KG_{\rm c}QN)$ entries for the coarse path matrices and
pilot responses, with $O(G_{\rm c}BM)$ temporary FFT storage.
